\section{绪论}\label{chap:introduction}

\subsection{研究背景与意义}

随着互联网技术的飞速发展和数字化进程的深入推进，网络已经成为人们日常生活中不可或缺的信息交流渠道。然而，网络通信的开放性和共享性也带来了严峻的信息安全挑战。政府、企业和个人的敏感信息在传输过程中面临着被窃取、篡改和监控的风险\supercite{fridrich2009steganography}。传统的加密技术虽然能够保护信息内容的机密性，但加密通信行为本身容易引起监控者的注意，从而招致进一步的审查和干预。

信息隐藏技术（Information Hiding）为解决上述问题提供了一种独特的思路。与加密技术不同，信息隐藏并不改变通信形式的外观，而是将秘密信息嵌入到看似正常的载体媒体（如图像、音频、视频等）中，使得通信双方的交互在外部观察者看来与正常通信无异\supercite{katzenbeisser2000information}。图像隐写术（Image Steganography）作为信息隐藏技术的重要分支，因其载体图像在网络中传输量大、格式多样、视觉冗余丰富等特点，成为研究最为活跃的方向之一。

早期的图像隐写方法主要基于空间域修改策略，如最低有效位（Least Significant Bit, LSB）替换\supercite{chan2004hiding}。这类方法实现简单，但嵌入操作会不可避免地改变载体图像的统计特性，容易被统计隐写分析方法检测\supercite{fridrich2012rich}。随着深度学习技术的发展，基于卷积神经网络的隐写分析工具（如 SRNet\supercite{boroumand2019deep}）展现出了极强的检测能力，对传统隐写方法构成了严重威胁。

为了从根本上解决隐写安全性问题，研究者提出了可证安全隐写（Provably Secure Steganography）的理论框架\supercite{hopper2009provably}。该理论指出，如果隐写系统生成的含密载体分布与正常载体分布在统计上不可区分（即 KL 散度为零），则任何检测器都无法以优于随机猜测的概率检测出隐写行为。近年来，基于深度生成模型（如变分自编码器 VAE、生成对抗网络 GAN、扩散模型 Diffusion Model）的隐写方案为实现可证安全提供了新的技术路径\supercite{de2022perfectly}。

与此同时，端到端加密（End-to-End Encryption, E2EE）技术的发展为即时通信安全提供了另一层保障。端到端加密确保只有通信双方能够解密消息内容，即使服务端被攻破也无法获取明文信息。Signal 协议\supercite{marlinspike2016signal}是目前最具代表性的端到端加密方案，被 WhatsApp、Signal 等主流通信应用广泛采用。将端到端加密与图像隐写相结合，可以实现"先加密、再隐写"的双重保护：即使隐写图像被提取出来，没有密钥也无法解密得到原始消息内容。

本文的研究正是在上述背景下展开的。本文选择基于扩散模型的 Pulsar 算法\supercite{jois2024pulsar}和 SparSamp 算法作为核心隐写方案，结合端到端加密机制，设计并实现一个完整的可证安全图像隐写通信系统。本文的工作目标包括以下两个方面：

\begin{enumerate}
    \item \textbf{算法工程化验证}：将 Pulsar 和 SparSamp 两种可证安全隐写算法集成到统一的服务框架中，验证其在真实通信环境下的嵌入容量、提取正确率和计算效率，并通过端到端加密与隐写的组合实现双重安全保护。
    \item \textbf{跨平台系统实现}：构建包含服务端（FastAPI）、Android 客户端（Jetpack Compose）和 Web 客户端（Vue.js 3）的完整隐写通信系统，实现从用户注册、好友管理到隐写消息收发的全流程功能，为可证安全隐写从理论走向工程应用提供可参考的实现方案。
\end{enumerate}

\subsection{国内外研究现状}

\subsubsection{传统图像隐写方法}

传统图像隐写方法可大致分为空间域方法和变换域方法两大类。

空间域隐写方法直接在图像像素值上进行修改。最经典的是 LSB 替换方法，通过将秘密信息的比特替换载体图像像素的最低有效位实现嵌入\supercite{chan2004hiding}。为了减小嵌入对图像统计特性的影响，研究者提出了 LSB 匹配\supercite{mielikainen2006lsb}、HUGO（Highly Undetectable Stego）\supercite{pevny2010using}、WOW（Wavelet Obtained Weights）\supercite{holub2012designing}和 S-UNIWARD\supercite{holub2014universal}等自适应隐写方法，通过最小化嵌入失真函数来优化嵌入位置的选择。

变换域隐写方法在图像的频域或其他变换域中嵌入信息。典型的方法包括基于离散余弦变换（DCT）的 JPEG 域隐写\supercite{westfeld2001f5}，如 F5 算法、nsF5 算法\supercite{fridrich2007statistically}和 J-UNIWARD\supercite{holub2014universal}等。

尽管上述方法在降低可检测性方面取得了显著进展，但由于嵌入操作本质上是对已有载体的修改，不可避免地会引入统计痕迹，无法在理论上保证绝对安全。

\subsubsection{基于深度学习的隐写分析}

隐写分析（Steganalysis）旨在检测载体中是否隐藏了秘密信息。传统隐写分析方法依赖手工设计的特征提取器，如 SRM（Spatial Rich Model）\supercite{fridrich2012rich}。随着深度学习的兴起，基于卷积神经网络的隐写分析方法展现了更优异的检测性能。代表性工作包括 Qian-Net\supercite{qian2015deep}、Xu-Net\supercite{xu2016structural}、Ye-Net\supercite{ye2017deep}和 SRNet\supercite{boroumand2019deep}等。这些方法能够自动学习隐写特征，对多种隐写方案都具有较高的检测准确率，进一步凸显了传统隐写方法在安全性方面的不足。

\subsubsection{可证安全隐写}

可证安全隐写的理论基础由 Hopper 等人\supercite{hopper2009provably}奠定。其核心思想是：隐写系统的安全性可以归结为含密载体分布与正常载体分布之间的统计距离（如 KL 散度）。若两者分布完全一致，则隐写系统是完美安全的（Perfectly Secure）。Cachin\supercite{cachin2004information}进一步从信息论角度形式化了这一安全准则。

实现可证安全的关键在于隐写编码过程不能对载体分布引入任何偏差。基于深度生成模型的隐写方案通过在模型的采样过程中嵌入信息，天然地满足了这一要求。代表性工作包括：

\begin{itemize}
    \item 基于 GAN 的方案：利用 GAN 的随机输入向量（latent code）嵌入信息\supercite{hu2018novel}。
    \item 基于 VAE 的方案：在变分自编码器的隐变量采样过程中嵌入信息。
    \item 基于 Flow 的方案：利用可逆归一化流模型的采样过程嵌入信息\supercite{de2022perfectly}。
    \item 基于扩散模型的方案：利用扩散模型逆向去噪过程中方差噪声的随机性嵌入信息。Pulsar\supercite{jois2024pulsar}是该方向的代表性工作，由 Jois 等人发表于 ACM CCS 2024，通过控制最后一步调度器的噪声来源实现逐比特编码，结合变码率纠错码在 256$\times$256 分辨率下平均嵌入 320--613 字节。SparSamp\supercite{wang2025sparsamp}则基于稀疏采样思想，通过消息驱动采样将消息段与伪随机数结合生成 MRN，在 IDDPM 的 P2-weighting 框架下逐像素采样嵌入消息，并通过逐步缩小候选消息集合消除解码歧义，单图容量可达数十千字节。
\end{itemize}

与其他生成模型相比，扩散模型生成的图像质量更高、多样性更好，是当前最有前景的可证安全隐写载体。

\subsubsection{端到端加密技术}

端到端加密（E2EE）是保障即时通信安全的核心技术。其基本思想是：消息在发送方设备上加密，只有接收方设备能够解密，中间节点（包括服务端）只能转发密文而无法获取明文。

Signal 协议\supercite{marlinspike2016signal}是目前最成熟的端到端加密方案，采用双棘轮算法（Double Ratchet Algorithm）实现前向安全和未来安全性。然而，Signal 协议的完整实现复杂度较高，需要维护大量的密钥状态和会话信息。

对于特定应用场景，可以采用更轻量的对称加密方案。通过预共享的助记词短语派生用户密钥，再基于通信双方的用户密钥协商出对称消息密钥，使用 AES-GCM 认证加密算法保护消息内容。这种方案在安全性上虽不如 Signal 协议完备（缺乏前向安全），但实现简单、性能优异，适合资源受限的移动端场景。

\subsection{主要研究内容}

本文围绕可证安全图像隐写系统的设计与实现，开展了以下主要工作：

\begin{enumerate}
    \item \textbf{双算法隐写引擎}：集成 Pulsar 和 SparSamp 两种基于扩散模型的可证安全隐写算法，设计了统一的算法注册与切换机制。Pulsar 基于 DDPM 框架，SparSamp 基于 IDDPM P2-weighting 框架，两者共享相同的 API 接口但底层实现独立。
    \item \textbf{端到端加密机制}：实现了基于对称助记词的端到端加密方案。用户通过助记词短语经 PBKDF2-SHA256 派生 256 位用户密钥；通信双方的用户密钥经 HKDF-SHA256 协商出对称消息密钥；消息使用 AES-256-GCM 加密封装为 SealedMessage 格式。
    \item \textbf{系统架构设计}：设计了基于客户端--服务端的跨平台隐写通信系统架构，服务端负责隐写算法计算和消息路由，客户端（Android 和 Web）负责用户交互、本地数据管理和加解密操作。
    \item \textbf{实时通信系统}：基于 WebSocket 协议实现了服务端与客户端之间的实时双向通信，支持消息即时推送、消息状态回执（已发送/已送达/已读）、消息撤回、在线状态感知和离线消息缓存。
    \item \textbf{用户体系与好友管理}：实现了完整的用户注册登录系统（JWT 认证）、基于邀请码的好友添加功能、好友请求的异步处理流程，以及单端登录安全策略。
    \item \textbf{多平台客户端开发}：分别使用 Kotlin Jetpack Compose 和 Vue.js 3 开发了功能完整的 Android 和 Web 客户端。Android 端采用 Material 3 Expressive 薄荷色主题，Web 端采用 Aegis Slate 暗色 Material 3 设计系统。两个客户端均实现了普通消息与隐写消息的无缝切换。
    \item \textbf{系统测试与验证}：对系统各项功能进行了全面测试，包括端到端加密正确性、双算法隐写嵌入提取、消息生命周期管理、跨平台兼容性等。
\end{enumerate}

\subsection{论文组织结构}

本文共分为六章，各章内容安排如下：

\textbf{第一章 绪论}：介绍了研究背景与意义，综述了国内外图像隐写术和端到端加密的研究现状，阐述了本文的主要研究内容和论文组织结构。

\textbf{第二章 相关技术与理论基础}：介绍了图像隐写术的基本概念、可证安全隐写的理论基础、Pulsar 和 SparSamp 两种隐写算法的工作原理、端到端加密的密钥协商与加密封装机制，以及系统开发所涉及的关键技术栈。

\textbf{第三章 系统设计}：从需求分析出发，详细阐述了系统的总体架构设计、服务端设计（API 接口、数据库、WebSocket 通信协议、离线消息队列、双算法隐写服务）、客户端设计（架构、本地存储、UI 导航）和安全设计（JWT 认证、密码安全、端到端加密、隐写密钥派生、单端登录）。

\textbf{第四章 系统实现}：分别介绍了服务端、Android 客户端和 Web 客户端的具体实现过程，包括项目结构、关键模块的类与方法设计，以及端到端加密封装、隐写消息收发、消息状态管理等核心功能的实现细节。

\textbf{第五章 系统测试与分析}：描述了系统的测试环境、功能测试（含端到端加密测试和双算法测试）、隐写性能分析和跨平台兼容性测试的过程与结果。

\textbf{第六章 总结与展望}：总结了本文的主要工作和贡献，分析了系统的不足之处，并对未来的研究方向进行了展望。
